\documentclass[11pt]{article}
\usepackage[margin=1in]{geometry}
\usepackage{amsmath,amssymb,amsthm}
\usepackage{graphicx}
\usepackage[hidelinks]{hyperref}

\newtheorem{theorem}{Theorem}
\newtheorem{corollary}{Corollary}
\newtheorem{remark}{Remark}

\newcommand{\N}{\omega}
\newcommand{\norm}[1]{\left\lVert #1\right\rVert}
\newcommand{\calP}{\mathcal P}
\newcommand{\calI}{\mathcal I}
\newcommand{\one}{\mathbf 1}
\newcommand{\eps}{\varepsilon}

\title{A Dilution Counterexample to the Quantitative \(\varepsilon\)-Strong \(Q\)-Measure Question}
\author{Agent lane 17}
\date{2026-06-16}

\begin{document}
\maketitle

\section*{Status}

This packet gives a strengthened conditional counterexample to Question 2 of
Aviles--Martinez-Cervantes--Poveda--Saenz \cite{AMPS2025}.  The earlier
version assumed a \(Q\)-point ultrafilter.  The present version only needs a
normalized \(\eta\)-strong \(Q\)-measure and then gives counterexamples for
every \(0<\varepsilon<\eta/2\).  In particular, if a strong \(Q\)-measure
exists, Question 2 has a negative answer for every \(0<\varepsilon<1/2\).

The packet also records the exact ZFC reduction: any \(\varepsilon\)-strong
measure whose mass on a bounded-per-block ideal is \(<1/2\) already yields the
same failure.  The remaining limitation is the same-parameter problem for an
arbitrary \(\varepsilon\)-strong measure.

\section*{Source Question}

The source paper asks the following question in Section 5.

\begin{center}
\includegraphics[width=0.95\linewidth]{figures/open_problem_crop.png}
\end{center}

In text: if \(\mu\) is an \(\varepsilon\)-strong \(Q\)-measure and
\((x_n)_{n\in\omega}\) is an \(L\)-orthogonal sequence in a Banach space \(X\),
must
\[
  \norm{x+\mu\hbox{-}\lim(x_n)}\geq \norm{x}+
  \varepsilon\norm{\mu\hbox{-}\lim(x_n)}
\]
hold for every \(x\in X\)?

\section*{Definitions Used}

For \(\varepsilon>0\), a finitely additive positive measure
\(\mu:\mathcal P(\omega)\to[0,1]\), vanishing on finite sets, is an
\(\varepsilon\)-strong \(Q\)-measure if every partition of \(\omega\) into
finite sets has a selector \(A\) with \(\mu(A)\ge\varepsilon\)
\cite[Definition 2.1]{AMPS2025}.

For a partition \(\mathcal P=(P_m)\) of \(\omega\) into finite sets, write
\[
  \mathcal I_{\mathcal P}=
  \left\{A\subseteq\omega:\sup_m |A\cap P_m|<\infty\right\}.
\]
The source paper defines a compact Cantor space \(K_{\mathcal P}\) and
coordinate functions \(f_n\in C(K_{\mathcal P})\), and recalls that
\((f_n)\) is \(L\)-orthogonal \cite[Definition 4.4 and Notation 4.5]{AMPS2025}.
It also proves the following estimate \cite[Lemma 4.6]{AMPS2025}: if \(\nu\)
is a positive finitely additive measure and \(\nu(I)\le \delta\norm{\nu}\) for
every \(I\in\mathcal I_{\mathcal P}\), then
\[
  \norm{\nu\hbox{-}\lim f_n-\norm{\nu}\one}\le 2\delta\norm{\nu}.
\]

For comparison with the source hierarchy, \(\mu\) is \(\delta\)-fit if for
every finite partition there is \(A\in\mathcal I_{\mathcal P}\) with
\(\mu(A)\ge\delta\mu(\omega)\), and it is fit if this holds for all
\(\delta<1/2\) \cite[Definition 4.7]{AMPS2025}.  Thus the first theorem below
is exactly the quantitative non-fit mechanism behind the counterexample.

\section*{Conditional Dependencies}

The ZFC reduction theorem has no extra set-theoretic dependency.  The relative
counterexample theorem assumes the existence of a normalized \(\eta\)-strong
\(Q\)-measure.  This assumption is weaker than requiring a \(Q\)-point when
\(\eta=1\): a \(Q\)-point gives a strong \(Q\)-measure, but the source paper
also shows under \(\mathfrak d=\operatorname{cov}(\mathcal M)\) that atomless
strong \(Q\)-measures exist \cite[Theorem 3.1]{AMPS2025}.

There is a genuine ZFC barrier to proving existence of such measures outright.
The source proves that it is consistent with ZFC that there are no
\(Q^+(\omega)\)-measures \cite[Theorem 1.2]{AMPS2025}; more explicitly,
Corollary 3.15 states that in Laver's model, the Mathias model, and Miller's
model there are no measure \(Q^+(\omega)\)-ideals.  Since every
\(\varepsilon\)-strong \(Q\)-measure is automatically a \(Q^+(\omega)\)-measure,
one cannot get a ZFC construction of an \(\varepsilon\)-strong measure that
works in all models.

\section*{Proof Intuition}

The source estimate says that in the canonical \(C(K_{\mathcal P})\) test
space, a measure whose bounded-per-block sets have mass \(<1/2\) has a
\(\mu\)-limit close to the constant function \(1\).  Then \(x=-1\) makes the
left side small, while the proposed quantitative lower bound is still larger
than \(1\).

The dilution trick manufactures exactly this situation from any sufficiently
strong source measure.  Put a small coefficient \(t<1/2\) on an
\(\eta\)-strong measure \(\sigma\), and put the remaining mass on a
block-density measure \(\lambda\) that vanishes on every
\(\mathcal I_{\mathcal P}\)-set.  The \(t\sigma\) part preserves enough
selector mass to be \(\varepsilon\)-strong whenever \(\varepsilon<t\eta\), but
the bounded-per-block ideal sees only at most \(t\) mass.

\begin{theorem}[ZFC non-fit reduction]
Let \(0<\varepsilon\), and let \(\mu:\mathcal P(\omega)\to[0,1]\) be a
positive finitely additive measure vanishing on finite sets.  Suppose that
\(\mu\) is an \(\varepsilon\)-strong \(Q\)-measure and that there are a
partition \(\mathcal P=(P_m)\) of \(\omega\) into finite sets and a scalar
\(\delta<1/2\) such that
\[
  \mu(I)\le \delta\mu(\omega)\qquad(I\in\mathcal I_{\mathcal P}).
\]
Then Question 2 fails for \(\mu\): there are a Banach space \(X\), an
\(L\)-orthogonal sequence \((x_n)\subseteq X\), and \(x\in X\) such that
\[
  \norm{x+\mu\hbox{-}\lim x_n}<
  \norm{x}+\varepsilon\norm{\mu\hbox{-}\lim x_n}.
\]
\end{theorem}

\begin{proof}
Set \(M=\mu(\omega)\).  Since \(\mu\) is \(\varepsilon\)-strong, \(M\ge
\varepsilon>0\).  Let \(X=C(K_{\mathcal P})\), and let \(x_n=f_n\) be the
coordinate functions from the source construction.  The sequence \((f_n)\) is
\(L\)-orthogonal.  By the source estimate,
\[
  \norm{\mu\hbox{-}\lim f_n-M\one}\le 2\delta M.
\]
Write \(u=\mu\hbox{-}\lim f_n\).  Since \(2\delta<1\),
\[
  \norm{u}\ge M-\norm{u-M\one}\ge M(1-2\delta)>0.
\]
Take \(x=-M\one\in X\).  Then
\[
  \norm{x+u}=\norm{u-M\one}\le 2\delta M<M,
\]
whereas
\[
  \norm{x}+\varepsilon\norm{u}
  \ge M+\varepsilon M(1-2\delta)>M.
\]
Therefore the proposed inequality is false.
\end{proof}

\begin{theorem}[Dilution counterexample]
Assume there is a normalized \(\eta\)-strong \(Q\)-measure
\(\sigma:\mathcal P(\omega)\to[0,1]\), so \(\sigma(\omega)=1\).  For every
\(0<\varepsilon<\eta/2\), there is an \(\varepsilon\)-strong \(Q\)-measure
\(\mu\) for which Question 2 fails.
\end{theorem}

\begin{proof}
Choose \(t\) with \(\varepsilon/\eta<t<1/2\).  Let
\(\mathcal P=(P_m)_{m\in\omega}\) be a partition of \(\omega\) into nonempty
finite sets with \(|P_m|\to\infty\).  Let \(\mathcal V\) be any free ultrafilter
on \(\omega\), and define
\[
  \lambda(A)=\lim_{m\to\mathcal V}\frac{|A\cap P_m|}{|P_m|}
  \qquad(A\subseteq\omega).
\]
Then \(\lambda\) is a positive finitely additive probability measure vanishing
on finite sets.  Also \(\lambda(I)=0\) for each \(I\in\mathcal I_{\mathcal P}\),
because \(|I\cap P_m|\) is bounded while \(|P_m|\to\infty\).

Define
\[
  \mu=t\sigma+(1-t)\lambda.
\]
This is again a positive finitely additive probability measure vanishing on
finite sets.

We first verify that \(\mu\) is \(\varepsilon\)-strong.  Let \(\mathcal R\) be
any partition of \(\omega\) into finite sets.  Since \(\sigma\) is
\(\eta\)-strong, there is a selector \(S\) for \(\mathcal R\) with
\(\sigma(S)\ge\eta\).  Hence
\[
  \mu(S)\ge t\sigma(S)\ge t\eta>\varepsilon.
\]
Thus \(\mu\) is \(\varepsilon\)-strong.

If \(I\in\mathcal I_{\mathcal P}\), then \(\lambda(I)=0\), so
\[
  \mu(I)=t\sigma(I)\le t.
\]
Choose \(\delta\) with \(t<\delta<1/2\).  The preceding theorem applies and
gives the desired failure of Question 2.
\end{proof}

\begin{corollary}
If a strong \(Q\)-measure exists, then Question 2 has a negative answer for
every \(0<\varepsilon<1/2\).  In particular, the previous Q-point
counterexample is the special case where the strong \(Q\)-measure is the
ultrafilter measure of a \(Q\)-point.
\end{corollary}

\begin{proof}
Normalize the strong \(Q\)-measure.  It is then \(1\)-strong, so the dilution
counterexample theorem applies with \(\eta=1\).
\end{proof}

\begin{remark}[What this does not close]
Starting from a particular \(\varepsilon\)-strong measure, the dilution argument
produces counterexamples only for smaller coefficients, namely any
\(\varepsilon'<\varepsilon/2\) after normalization.  A same-parameter
counterexample for every given \(\varepsilon\)-strong measure would require a
new argument showing that the measure itself, or a same-strength modification
of it, has the non-fit/block-smallness property above.  This is closely tied
to the source paper's Question 3, which asks whether the notions
\(Q^+(\omega)\)-measure, \(\delta\)-fit, fit, \(\varepsilon\)-strong, and
strong \(Q\)-measure are equivalent or distinct.
\end{remark}

\section*{Verification Notes}

The proof has two scalar requirements.  In the reduction theorem,
\(\delta<1/2\) gives
\[
  2\delta M<M<M+\varepsilon M(1-2\delta).
\]
In the dilution theorem, \(t<1/2\) gives the block-smallness constant, while
\(\varepsilon<t\eta\) preserves \(\varepsilon\)-strongness.  The included script
\[
\texttt{code/check\_scalar\_inequality.py}
\]
checks representative scalar choices; the mathematical proof above does not
depend on computation.

\section*{Limitations and Novelty Check}

This remains a conditional packet because it does not prove, in ZFC, the
existence of an \(\eta\)-strong \(Q\)-measure.  Such a proof is impossible in
the naive existence sense, since the source paper gives models with no
\(Q^+(\omega)\)-measures, and \(\varepsilon\)-strong implies \(Q^+(\omega)\).
The result also does not address coefficients \(\varepsilon\ge 1/2\), where
the present block-smallness contradiction cannot work.

I searched the run indexes for arXiv:2501.06537, \(L\)-orthogonal sequences and
elements, \(Q\)-measures, \(\varepsilon\)-strong \(Q\)-measures, fit
\(Q\)-measures, and the quantitative Theorem 4.3 question; only this packet
hit.  A bounded web search on 2026-06-16 for exact phrases including
``epsilon-strong Q-measure'', ``fit Q-measure'', ``2501.06537'', and
``quantitative version Theorem 4.3 L-orthogonal'' did not reveal a later answer
to Question 2.

\section*{Human Review Recommendation}

Review the dilution step \(\mu=t\sigma+(1-t)\lambda\), especially the two
inequalities \(\mu(S)\ge t\eta>\varepsilon\) for arbitrary finite partitions
and \(\mu(I)\le t<\delta<1/2\) for \(I\in\mathcal I_{\mathcal P}\).  The
remaining mathematical frontier is the same-parameter case for an arbitrary
\(\varepsilon\)-strong measure, not a gap in the proved relative result.

\begin{thebibliography}{9}
\bibitem{AMPS2025}
A. Aviles, G. Martinez-Cervantes, A. Poveda, and L. Saenz,
\emph{A Banach space with \(L\)-orthogonal sequences but without
\(L\)-orthogonal elements},
arXiv:2501.06537, 2025.
\end{thebibliography}

\end{document}
